\section*{第 \ref{ch:b2:hermitian} 章 --- 埃尔米特形式}

\begin{solution}{exo:b2:hermitian:1}
$\langle x, y\rangle = \conj{1}\cdot\iu + \conj{\iu}\cdot 1 = \iu -
\iu = 0$：正交。$\norm x = \norm y = \sqrt{1 + 1} = \sqrt2$。
标准正交基：$\bigl(\frac{1}{\sqrt2}(1, \iu),\;
\frac{1}{\sqrt2}(\iu, 1)\bigr)$ --- 即归一化后的这一对向量本身。
\end{solution}

\begin{solution}{exo:b2:hermitian:2}
第一个矩阵：与其共轭转置相等（对角元为实数，$\conj{\iu} =
-\iu$ 互为共轭）：埃尔米特矩阵（从而为正规矩阵）；不是酉矩阵（$A^\dagger A
\neq I$：各列不是单位向量）。

第二个矩阵：$A^\dagger A = \frac12\begin{pmatrix} 1 & -\iu\\ -\iu &
1\end{pmatrix}\begin{pmatrix} 1 & \iu\\ \iu & 1\end{pmatrix} =
\frac12\begin{pmatrix} 2 & 0\\ 0 & 2\end{pmatrix} = I$：酉矩阵
（从而为正规矩阵）；不是埃尔米特矩阵。

第三个矩阵：$A^\dagger A = E_{22} \neq E_{11} = AA^\dagger$：不是正规矩阵
（因此既不是埃尔米特矩阵，也不是酉矩阵） --- 经典的幂零
反例。
\end{solution}

\begin{solution}{exo:b2:hermitian:3}
唯一性：设 $A = H + \iu K$，其中 $H^\dagger = H$，$K^\dagger = K$，
则 $A^\dagger = H - \iu K$，于是 $H = \frac{A + A^\dagger}{2}$，
$K = \frac{A - A^\dagger}{2\iu}$；这两个矩阵均为埃尔米特矩阵
（验证：$\bigl(\frac{A - A^\dagger}{2\iu}\bigr)^\dagger = \frac{
A^\dagger - A}{-2\iu} = K$），且可还原 $A$，即存在性成立。

正规性：$A^\dagger A - AA^\dagger = (H - \iu K)(H + \iu K) - (H
+ \iu K)(H - \iu K) = 2\iu(HK - KH)$：该差为零当且仅当 $HK = KH$。
\end{solution}

\begin{solution}{exo:b2:hermitian:4}
$\chi_A = (X-2)^2 - 1$：特征值为 $3$ 与 $1$。特征向量由
直接计算得到：
\[
A\begin{pmatrix}1\\ -\iu\end{pmatrix}
= \begin{pmatrix} 2 + \iu(-\iu)\\ -\iu + 2(-\iu)\end{pmatrix}
= \begin{pmatrix} 3\\ -3\iu \end{pmatrix}
= 3\begin{pmatrix}1\\ -\iu\end{pmatrix},
\qquad
A\begin{pmatrix}1\\ \iu\end{pmatrix}
= \begin{pmatrix} 2 + \iu\cdot\iu\\ -\iu + 2\iu\end{pmatrix}
= \begin{pmatrix}1\\ \iu\end{pmatrix}.
\]
标准正交特征基：$u_1 = \frac{1}{\sqrt2}(1, -\iu)$
（对应特征值 $3$），$u_2 = \frac{1}{\sqrt2}(1, \iu)$（对应特征值
$1$）；正交性见\cref{exo:b2:hermitian:1}。求幂时利用
谱投影 $A^k = 3^k u_1u_1^\dagger + 1^k\,
u_2u_2^\dagger$：
\[
A^k = U\begin{pmatrix} 3^k & 0\\ 0 & 1\end{pmatrix}U^\dagger
= \frac{3^k}{2}\begin{pmatrix} 1 & \iu\\ -\iu & 1\end{pmatrix}
+ \frac{1}{2}\begin{pmatrix} 1 & -\iu\\ \iu & 1\end{pmatrix}
= \frac12\begin{pmatrix}
3^k + 1 & (3^k - 1)\iu\\
-(3^k-1)\iu & 3^k + 1
\end{pmatrix}.
\]
（验证 $k = 1$ 时即得 $A$。）
\end{solution}

\begin{solution}{exo:b2:hermitian:5}
紧：既是闭集（$I$ 在连续映射 $U \mapsto
U^\dagger U$ 下的原像），又有界（各列为单位向量，元素
模长 $\leq 1$），都在 $\mathcal{M}_n(\C) \simeq \R^{2n^2}$ 中。

特征值：$\operatorname{diag}(\lambda, 1, \dots, 1)$
对任意 $\abs\lambda = 1$ 均为酉矩阵。行列式：$\abs{\det U}^2 = \det
U^\dagger \det U = \det(U^\dagger U) = 1$（利用 $\det A^\dagger =
\conj{\det A}$）：$\det U$ 位于单位圆周上。
\end{solution}

\begin{solution}{exo:b2:hermitian:6}
由谱定理，在 $H$ 的标准正交特征基中考虑：
一切归结为标量 $\lambda \in \R$（即特征值）的计算。
$H + \iu I$ 的特征值为 $\lambda + \iu \neq 0$，故可逆。$U$
的特征值为 $\mu = \frac{\lambda - \iu}{\lambda + \iu}$，其模
为 $1$（$\abs{\lambda - \iu} = \abs{\lambda + \iu}$，当
$\lambda$ 为实数）：$U^\dagger U = I$ 成立，因 $U$ 可酉对角化，
其特征值的模均为 1（在所取的标准正交基中
为对角矩阵）。若 $\mu = 1$，则 $-\iu = \iu$，
不可能，故 $1 \notin \operatorname{Sp} U$。（凯莱变换把埃尔米特矩阵映为
不以 1 为特征值的酉矩阵 --- 即
从 $\R$ 到单位圆的映射的矩阵版本。）
\end{solution}

\begin{solution}{exo:b2:hermitian:7}
对特征对 $Ax = \lambda x$（$x \neq 0$）：$\lambda\norm x^2 =
\langle x, Ax\rangle > 0$，故 $\lambda > 0$（已
由\cref{prop:b2:hermitian:eigenvalues}知其为实数）。平方根：在谱基下
令 $B = U\operatorname{diag}(\sqrt{\lambda_i})U^\dagger$：
埃尔米特、正定，且 $B^2 = A$。行列式为
各正特征值之积。
\end{solution}

\begin{solution}{exo:b2:hermitian:8}
\begin{enumerate}
  \item $\norm{u(x)}^2 = \langle u(x), u(x)\rangle = \langle x,
        u^*u(x)\rangle = \langle x, uu^*(x)\rangle =
        \norm{u^*(x)}^2$。将此用于正规自同态 $u -
        \lambda\,\mathrm{id}$（其伴随为 $u^* -
        \conj\lambda$，且正规性可继承）：$\norm{(u -
        \lambda)x} = \norm{(u^* - \conj\lambda)x}$，故二者
        的核相同。
  \item 对特征向量 $u(x) = \lambda x$，$u(y) = \mu y$
        （$\lambda \neq \mu$）：由 (1)，$u^*(x) = \conj\lambda x$；
        于是
        \[
        \lambda\langle y, x\rangle = \langle y, u(x)\rangle
        = \langle u^*(y), x\rangle = \langle \conj\mu\, y, x\rangle
        = \mu \langle y, x\rangle ,
        \]
        故 $\langle y, x\rangle = 0$。$E_\lambda^\perp$ 的稳定性：
        对 $x \perp E_\lambda$ 及 $z \in E_\lambda$，$\langle z,
        u(x)\rangle = \langle u^*(z), x\rangle =
        \langle\conj\lambda z, x\rangle = 0$。
  \item 对维数归纳：在 $\C$ 上 $u$ 有一个特征向量
        $e_1$（单位化）；其正交补
        在 $u$ 下稳定（由 (2)）\emph{且}在 $u^*$ 下稳定（同理，
        角色互换），故限制后仍为正规自同态：对其归纳
        并拼接标准正交特征基。反之，
        酉可对角化的 $u = UDU^\dagger$ 满足 $u^*u
        = U\conj D D U^\dagger = U D\conj D U^\dagger = uu^*$：
        即为正规自同态。
\end{enumerate}
\end{solution}

\begin{solution}{exo:b2:hermitian:9}
特征值：对 $u(x) = \lambda x$，$x \neq 0$：
\[
\lambda\norm x^2 = \langle x, u(x)\rangle
= \langle u^*(x), x\rangle = -\langle u(x), x\rangle
= -\conj{\langle x, u(x)\rangle} = -\conj\lambda\,\norm x^2 ,
\]
于是 $\lambda = -\conj\lambda$，即为纯虚数。由于 $(\iu
u)^* = -\iu\,u^*$（伴随对标量是共轭线性的），
由 $u^* = u$ 得 $(\iu u)^* = -\iu u$：映射 $u \mapsto \iu u$
将埃尔米特映为反埃尔米特，其逆为 $w \mapsto -\iu
w$，是双射。反埃尔米特的 $u$ 满足 $u^*u = -u^2 =
uu^*$，故为正规的，从而
由\cref{exo:b2:hermitian:8}可知其可酉对角化。
\end{solution}

\begin{solution}{exo:b2:hermitian:10}
\begin{enumerate}
  \item $S$ 置换一组标准正交基：$\norm{Sx} = \norm
        x$，故 $S$ 为酉的。坐标按 $j = 0,
        \dots, n-1$ 模 $n$ 编号：$(Sx)_j = x_{j-1}$，于是对
        $(f_k)_j = \frac{\omega^{jk}}{\sqrt n}$：
        \[
        (Sf_k)_j = \frac{\omega^{(j-1)k}}{\sqrt n}
        = \omega^{-k}\,(f_k)_j :
        \qquad Sf_k = \omega^{-k}f_k .
        \]
        标准正交性：$\langle f_k, f_l\rangle = \frac1n
        \sum_j \omega^{j(l-k)} = \delta_{kl}$（非平凡单位根的等比和
        为零）。
  \item $Cf_k = \sum_j c_j S^jf_k = \bigl(\sum_j
        c_j\omega^{-jk}\bigr)f_k = \widehat c(k)\,f_k$：每个
        循环矩阵在标准正交傅里叶基下为对角的，
        从而正规，谱为 $\{\widehat c(k)\}$。（
        基变换即离散傅里叶变换：
        卷积化为乘积。）
\end{enumerate}
\end{solution}

\begin{solution}{exo:b2:hermitian:11}
（$\Leftarrow$）设 $P^2 = P = P^*$。每个 $v$ 可分解为 $v = Pv
+ (v - Pv)$，其中 $Pv \in \operatorname{im} P$，且 $P(v - Pv) =
0$。两部分彼此正交：对任意 $x, y$，
\[
\langle Px, (I - P)y\rangle = \langle x, P(I-P)y\rangle
= \langle x, (P - P^2)y\rangle = 0 :
\]
$\ker P \perp \operatorname{im} P$，故 $P$ 是
其像上的正交投影。（$\Rightarrow$）若 $P$ 为
到 $F = \operatorname{im}P$ 的正交投影：对所有
$x, y$，$\langle Px, y\rangle = \langle Px, Py\rangle$（
分量 $y - Py \perp F$ 消去），且对称地有 $\langle
x, Py\rangle = \langle Px, Py\rangle$：$\langle Px, y\rangle =
\langle x, Py\rangle$，即 $P^* = P$。矩阵形式：投影到 $\C v$
（$\norm v = 1$）：$Px = v\,\langle v, x\rangle$，即 $P =
vv^\dagger$；投影到 $\operatorname{Vect}(v_1, \dots, v_k)$ 上，
标准正交时：$P = \sum_i v_iv_i^\dagger$。
\end{solution}

\begin{solution}{exo:b2:hermitian:12}
对 $A = U D U^\dagger$ 作对角化（谱定理），$D$ 为对角阵，
其对角元取自 $\lambda_i$。于是 $P_i = L_i(A) = U
L_i(D)U^\dagger$，且 $L_i(D)$ 为对角阵，其对角元为
$L_i(\lambda_j) = \delta_{ij}$：恰在
$E_i$ 对应的位置处为 1。于是 $P_i$ 为埃尔米特的（$L_i$ 与 $D$ 皆为实的），幂等，
其像为 $E_i$，核为 $\bigoplus_{j\neq i}E_j = E_i^\perp$
（特征子空间两两正交），即到
$E_i$ 上的正交投影（见\cref{exo:b2:hermitian:11}）。不相交的对角模式
给出 $P_iP_j = 0$（$i \neq j$）；$\sum_i L_i = 1$（次数 $< p$，
值为 $1$ 的点有 $p$ 个），故 $\sum P_i = I$；且 $\sum_i
\lambda_iL_i(\lambda_j) = \lambda_j$ 给出 $A = \sum
\lambda_iP_i$。对任意多项式 $f$：$f(D)$ 的对角元为
$f(\lambda_j)$，于是
\[
f(A) = \sum_{i=1}^{p} f(\lambda_i)\,P_i :
\]
$A$ 的函数借助谱来计算 --- 此演算
本科三年级分册将其推广到连续 $f$ 乃至更一般的情形。
\end{solution}

\begin{solution}{pb:b2:hermitian:1}
\textbf{1.} $\conj{\langle x, Ax\rangle} = \langle Ax, x\rangle
= \langle x, A^*x\rangle = \langle x, Ax\rangle$：为实数。记
$x = \sum c_ie_i$：
\[
R_A(x) = \frac{\sum_i\lambda_i\abs{c_i}^2}
{\sum_i\abs{c_i}^2} ,
\]
为特征值的加权平均：它位于
$\intcc{\lambda_n}{\lambda_1}$ 内，其界在
$e_1$ 与 $e_n$ 处取到。

\textbf{2.} 对实数 $t$ 及任意 $v$，展开 $R_A(x + tv) =
\frac{N(t)}{D(t)}$，其中
\[
N(t) = \langle x, Ax\rangle + 2t\Re\langle v, Ax\rangle +
t^2\langle v, Av\rangle,
\quad
D(t) = \norm x^2 + 2t\Re\langle v, x\rangle + t^2\norm v^2 .
\]
在 $t = 0$ 处的导数为
\[
\frac{2}{\norm x^2}\,
\Re\bigl\langle v,\ Ax - R_A(x)\,x\bigr\rangle .
\]
该导数对所有 $v$ 为零，当且仅当 $\Re\langle v, w\rangle = 0$ 对
所有 $v$ 成立，其中 $w = Ax - R_A(x)x$；将 $v$ 换为 $\iu v$
同样消去虚部：$w = 0$，即 $Ax = R_A(x)x$。
$R_A$ 的临界点恰为特征向量，
其临界值即为相应的特征值。

\textbf{3.} 设 $x = \sum_{i\leq k}c_ie_i \in V_k$，则 $R_A(x)$ 是
$\lambda_1, \dots, \lambda_k$ 的加权平均，因此
$\geq \lambda_k$，在 $e_k$ 处取等号：$\min_{V_k} R_A =
\lambda_k$。类似地，在 $W_k$ 上，该平均涉及
$\lambda_k, \dots, \lambda_n$：$\max_{W_k}R_A = \lambda_k$。

\textbf{4.} 设 $\dim V = k$。则 $\dim V + \dim W_k = n + 1 >
n$，故存在单位向量 $x \in V \cap W_k$，且 $R_A(x) \leq
\lambda_k$（第 3 题）：$\min_{V}R_A \leq \lambda_k$ 对每个
这样的 $V$ 成立。由于 $V_k$ 取到 $\lambda_k$，故极大–极小值等于
$\lambda_k$。极小–极大公式的论证与此相同，只是角色对调
（由 $\dim W = n - k + 1$ 可知 $W \cap V_k \neq
\{0\}$，故 $\max_W R_A \geq \lambda_k$，且在 $W_k$ 处取到）。

\textbf{5.} $R_B(x) = R_A(x) + \frac{\langle x,
(B-A)x\rangle}{\norm x^2} \geq R_A(x)$ 逐点成立。
取 $\min$（对任意 $k$ 维的 $V$），再取 $\max$（对 $V$）：
$\lambda_k(B) \geq \lambda_k(A)$，由第 4 题可得。

\textbf{6.} 由格拉斯曼公式：$\dim(V\cap W) = \dim V + \dim W -
\dim(V + W) \geq \dim V + \dim W - n > 0$。两次应用此式：
\[
\dim(V_1\cap V_2\cap V_3)
\geq \dim(V_1\cap V_2) + \dim V_3 - n
\geq \dim V_1 + \dim V_2 + \dim V_3 - 2n .
\]

\textbf{7.} 子空间 $W_i(A)$、$W_j(B)$（第 3 题，针对
$A$ 与 $B$），以及 $V_{i+j-1}(A+B)$ 的维数分别为 $(n-i+1) +
(n-j+1) + (i+j-1) = 2n + 1 > 2n$：由第 6 题可知存在单位向量
$x$ 同时属于这三者。于是
\[
\lambda_{i+j-1}(A+B) \leq R_{A+B}(x)
= R_A(x) + R_B(x) \leq \lambda_i(A) + \lambda_j(B),
\]
左侧不等式成立，因为 $x \in V_{i+j-1}(A+B)$（第
3 题），右侧则由两个 $W$ 给出。

\textbf{8.} 在 $E$ 的谱基下：$\norm{Ex}^2 = \sum
\lambda_k(E)^2\abs{c_k}^2 \leq \max_k\lambda_k(E)^2\,\norm
x^2$，在对应的特征向量处取到：$\vertiii E_2
= \max_k\abs{\lambda_k(E)}$。由外尔扰动不等式（取 $j = 1$）：$\lambda_k(A+E)
\leq \lambda_k(A) + \lambda_1(E) \leq \lambda_k(A) + \vertiii
E_2$。将此式用于 $(A+E) + (-E)$，得 $\lambda_k(A) \leq
\lambda_k(A+E) + \vertiii E_2$。综合两式：
$\abs{\lambda_k(A+E) - \lambda_k(A)} \leq \vertiii E_2$。

\textbf{9.} 下界由 $P \geq 0$ 及第 5 题得出。上界：$P$
秩为 $1$，故 $\lambda_2(P) = 0$；由外尔扰动不等式，取 $i = k-1$，$j =
2$:
\[
\lambda_k(A + P) \leq \lambda_{k-1}(A) + \lambda_2(P)
= \lambda_{k-1}(A) .
\]

\textbf{10.} $A + E = \begin{pmatrix} 2 & 1+\iu\\ 1-\iu &
2\end{pmatrix}$：特征多项式为 $(2-\lambda)^2 -
\abs{1+\iu}^2 = (2-\lambda)^2 - 2$，谱为 $\{2 + \sqrt2,\ 2
- \sqrt2\}$。与 $\operatorname{Sp}A = \{3, 1\}$ 相比：
\[
\abs{(2+\sqrt2) - 3} = \abs{(2-\sqrt2) - 1} = \sqrt2 - 1
\approx 0.414 \leq 1 = \vertiii E_2 . \checkmark
\]

\textbf{11.} 考虑 $\C^{n-1} = \operatorname{Vect}(e_1, \dots,
e_{n-1})$ 嵌入 $\C^n$ 中（取标准基）：对其中的 $x$，
$\langle x, Bx\rangle = \langle x, Ax\rangle$，故 $R_B$ 是
$R_A$ 在该子空间上的限制。上界：$\lambda_k
(B)$ 的极大–极小取遍 $k$ 维子空间，这些子空间\emph{位于}
$\C^{n-1}$ 中，它们是 $\C^n$ 中子空间族的一个子族：$\lambda_k(B) \leq
\lambda_k(A)$。下界：对 $\lambda_k(B)$
取遍 $\C^{n-1}$ 的子空间，维数为 $(n-1)-k+1 =
n-k$。其中每个也是 $\C^n$ 中维数为 $n -
(k+1) + 1$，故其极大值 $\geq \lambda_{k+1}(A)$：
$\lambda_k(B) \geq \lambda_{k+1}(A)$。

\textbf{12.} 逐次删去行与列，并反复运用
第 11 题：每删去一行或一列，下标便下移一位，由此得
$\lambda_{k+m}(A) \leq \lambda_k(B) \leq \lambda_k(A)$。

\textbf{13.} 以置换矩阵（酉矩阵）对 $A$ 作共轭，
既不改变谱，也不改变对角元的
多重集。不妨设 $d_1, \dots, d_k$ 占据前面的位置。
于是前 $k\times k$ 阶主子矩阵 $B$ 有
$\operatorname{tr} B = \sum_{i\leq k}d_i$，其特征值
满足 $\mu_i(B) \leq \lambda_i(A)$（第 12 题）；相加得
$\sum_{i\leq k}d_i \leq \sum_{i\leq k}\lambda_i(A)$。当 $k =
n$ 时，两边均为 $\operatorname{tr} A$。

\textbf{14.} 取 $x_i = e_i$，所得值为
$\sum_{i\leq k}\lambda_i$，故极大值 $\geq$ 该值。反之，任一
标准正交族 $(x_1, \dots, x_k)$ 可扩充为一组
标准正交基，即得酉矩阵 $U$，其前若干列为
$x_i$；则 $\sum_i\langle x_i, Ax_i\rangle$ 为
前 $k$ 个对角元（取自 $U^\dagger AU$）之和，由第
13 题知其至多为其 $k$ 个最大特征值之和 ---
即 $\sum_{i\leq k}\lambda_i(A)$。范凯
极大原理由此得出。

\textbf{15.} 交错性（$\operatorname{Sp}A = \{2+\sqrt2, 2,
2-\sqrt2\}$，$\operatorname{Sp}B = \{3, 1\}$）：
\[
2 \leq 3 \leq 2 + \sqrt2,
\qquad
2 - \sqrt2 \leq 1 \leq 2 . \checkmark
\]
由舒尔定理（对角元为 $(2,2,2)$）：$2 \leq 2+\sqrt2$；$4 \leq 4 +
\sqrt2$；$6 = 6$（迹）。\checkmark

\textbf{16.} $b_{ii} - a_{ii} = \langle e_i, (B-A)e_i\rangle
\geq 0$；求和即得迹的不等式。对任意 $C$：$\langle x,
C^\dagger(B - A)Cx\rangle = \langle Cx, (B-A)(Cx)\rangle \geq
0$：$C^\dagger AC \leq C^\dagger BC$。

\textbf{17.} $A \geq 0$ 显然；$B - A = \begin{pmatrix} 1 & 1\\
1 & 1\end{pmatrix}$ 为半正定（特征值 $2,
0$）：$A \leq B$。但
\[
B^2 = \begin{pmatrix} 5 & 3\\ 3 & 2\end{pmatrix},
\qquad
B^2 - A^2 = \begin{pmatrix} 4 & 3\\ 3 & 2\end{pmatrix},
\qquad \det(B^2 - A^2) = -1 < 0 :
\]
不是半正定的。平方运算不保持
勒夫纳序。

\textbf{18.} 设 $S = \sqrt A$，$T = \sqrt B$（埃尔米特
半正定矩阵，将\cref{exo:b2:hermitian:7}推广到
半正定情形，用同样的谱公式）。$T - S$ 为
埃尔米特矩阵；设 $\mu$ 为任一特征值，$v$ 为单位特征向量。
由 $T^2 - S^2 = T(T - S) + (T - S)S$，有
\[
0 \leq \langle v, (B - A)v\rangle
= \langle Tv, (T-S)v\rangle + \langle (T-S)v, Sv\rangle
= \mu\bigl(\langle v, Tv\rangle + \langle v,
Sv\rangle\bigr)
\]
（$\mu$ 为实数）。若 $\langle v, Tv\rangle + \langle v,
Sv\rangle > 0$，则 $\mu \geq 0$。若其为零，则两个
非负项皆为零；$\langle v, Tv\rangle =
\norm{T^{1/2}v}^2 = 0$ 推出 $Tv = 0$，同理 $Sv = 0$，于是
$\mu v = (T - S)v = 0$，从而 $\mu = 0$。$T -
S$ 的全部特征值皆 $\geq 0$，即 $\sqrt A \leq \sqrt B$。

\textbf{19.} 对 $A^{-1/2}$ 作合同变换（第 16 题）：$I \leq M
:= A^{-1/2}BA^{-1/2}$。于是 $M$ 的全部特征值皆 $\geq 1$，
从而 $M^{-1}$ 的特征值落在 $\intoc{0}{1}$ 中，故 $M^{-1} \leq I$。
但 $M^{-1} = A^{1/2}B^{-1}A^{1/2}$。将 $M^{-1} \leq
I$ 以 $A^{-1/2}$ 作合同变换，得 $B^{-1} \leq A^{-1}$。

\textbf{20.} $\sqrt A^{-1}(AB)\sqrt A = \sqrt A\,B\sqrt A$，故
$AB$ 相似于埃尔米特正定矩阵 $\sqrt
A\,B\sqrt A$（正定性：$\langle x, \sqrt AB\sqrt Ax\rangle =
\langle \sqrt Ax, B\sqrt Ax\rangle > 0$）。因此其特征值
皆为实数且为正。此外
\[
\langle x, \sqrt AB\sqrt Ax\rangle
\leq \lambda_{\max}(B)\,\norm{\sqrt Ax}^2
= \lambda_{\max}(B)\,\langle x, Ax\rangle
\leq \lambda_{\max}(A)\lambda_{\max}(B)\norm x^2 ,
\]
于是 $\lambda_{\max}(AB) = \max R_{\sqrt AB\sqrt A} \leq
\lambda_{\max}(A)\lambda_{\max}(B)$。

\textbf{21.} 取 $v_k = (\sin j\theta_k)_j$，并利用恒等式
$\sin((j-1)\theta) + \sin((j+1)\theta) =
2\sin(j\theta)\cos\theta$，对 $2 \leq j \leq n-1$，
\[
(T_nv_k)_j = -\sin((j{-}1)\theta_k) + 2\sin(j\theta_k) -
\sin((j{+}1)\theta_k)
= (2 - 2\cos\theta_k)\sin(j\theta_k) .
\]
第 $1$ 行成立，因为 $\sin(0\cdot\theta_k) = 0$，第 $n$
行成立，因为 $\sin((n+1)\theta_k) = \sin(k\pi) = 0$。边界
条件恰好选出 $\theta_k = \frac{k\pi}{n+1}$。于是
$T_nv_k = 4\sin^2\bigl(\frac{k\pi}{2(n+1)}\bigr)v_k$；这 $n$
个值在 $\intoo04$ 中互不相同，且 $v_k \neq 0$：这便是
全部谱，皆为正，故 $T_n$ 为正定矩阵。

\textbf{22.} $\min_id_i\,I \leq D \leq \max_id_i\,I$，故 $T_n +
\min_id_i\,I \leq T_n + D \leq T_n + \max_id_i\,I$
（加 $T_n$ 保持序关系）；第 5 题及
$\lambda_k(T_n + cI) = \lambda_k(T_n) + c$ 给出此夹逼。

\textbf{23.} $T_2 = \begin{pmatrix} 2 & -1\\ -1 &
2\end{pmatrix}$ 的谱为 $\{3, 1\}$，且
\[
2 - \sqrt2 \;\leq\; 1 \;\leq\; 2 \;\leq\; 3 \;\leq\; 2 +
\sqrt2 :
\]
柯西交错性已验证。（这些谱与
第 15 题中的相同：以 $\operatorname{diag}(1,-1,1)$ 共轭会
使非对角元变号。）图论解读：$T_2$ 是
去掉末顶点的路径的类拉普拉斯矩阵，
--- 即一个主子阵，恰为
第 11 题的情形。

\textbf{24.} 最大与最小特征值的渐近行为如下：
\[
\lambda_{\min}(T_n) = 4\sin^2\frac{\pi}{2(n+1)}
\sim \frac{\pi^2}{(n+1)^2},
\qquad
\lambda_{\max}(T_n) = 4\cos^2\frac{\pi}{2(n+1)}
\longrightarrow 4 .
\]
由此
\[
\kappa_n = \frac{\lambda_{\max}}{\lambda_{\min}}
\sim \frac{4(n+1)^2}{\pi^2} :
\]
网格越细，离散二阶
导数的条件越差 --- 这一事实推动了数值
线性代数方法的设计。

\textbf{25.} (i) 特征多项式的根在一般矩阵的扰动下
可能剧烈移动，但极小极大刻画
把每个埃尔米特特征值夹在
显式的优化值之间，迫使 $1$-利普希茨
稳定性（第 8--9 题）。(ii) 第 1、5、16--20 题仅用到
瑞利商的极值；外尔、交错、舒尔与
凯·范（第 7--14 题）确实需要完整的
子空间极小极大。(iii) 勒夫纳序把这些标量
不等式化为矩阵不等式的演算 --- 其中
既有真正的陷阱（第 17 题），也有真正的定理（第 18--19 题）。
(iv) 这些工具是数值分析的家常便饭
（第 24 题的刚度矩阵）、量子微扰
理论（外尔：能级移动至多为扰动的范数
），以及本科三年级分册中关于
紧自伴算子的极小极大原理。
\end{solution}
